\documentclass[11pt]{article}
\usepackage{../../latex/stat4ml_style}

\title{Chapter 7: Bandits and Policy-Gradient Reinforcement Learning \\ \large Part III --- Statistical Foundations of Deep Learning \& Reinforcement Learning}
\author{Omid Safarzadeh}
\date{}

\begin{document}
\maketitle

\begin{abstract}
\noindent
Two statistical problems sit under reinforcement learning's action side.
The first is the explore--exploit trade-off, in its purest form the
\keyterm{multi-armed bandit}: pull arms to maximise reward while learning which
arm is best, and the right amount of exploration falls out of a concentration
inequality (UCB) or a posterior (Thompson sampling), with regret that grows only
logarithmically. The second is \emph{optimising} a policy you can only sample
from --- and its answer is the \keyterm{policy gradient}, which is nothing but the
score function of Chapter~2 applied to a policy. The gradient of expected reward
equals the expected reward times the score, an identity that lets you
differentiate through a sampling step, and subtracting a baseline --- a control
variate --- cuts its variance without biasing it.
\end{abstract}

\tableofcontents

% ======================================================================
\section{Bandits: exploration from a confidence bound}
% ======================================================================

A \keyterm{$K$-armed bandit} offers arms with unknown mean rewards
$\mu_1,\dots,\mu_K$; each round you pull one and see a noisy reward. The cost of
not knowing the best arm $\mu^\star=\max_k\mu_k$ is the \keyterm{regret}
$\mathrm{Reg}(T)=\sum_{t=1}^T(\mu^\star-\mu_{a_t})$. Pulling greedily can lock
onto a wrong arm forever (linear regret); exploring at a constant rate
($\epsilon$-greedy) wastes a constant fraction of pulls, also linear. The art is
to explore \emph{just enough}, and Hoeffding says how much.

\begin{definition}[Upper confidence bound]
After $n_k$ pulls of arm $k$ with empirical mean $\hat\mu_k$, the \keyterm{UCB}
index is
\[
  \hat\mu_k + \sqrt{\frac{c\log t}{n_k}},
\]
and UCB1 pulls the arm of largest index.
\end{definition}

The bonus is exactly a Hoeffding confidence width (Chapter~1): with high
probability $\mu_k\le\hat\mu_k+\sqrt{c\log t/n_k}$, so the index is an
\emph{optimistic} estimate of the arm's value. Optimism forces exploration of
under-sampled arms (small $n_k$, large bonus) but lets the bonus shrink as
evidence accumulates, which yields \keyterm{logarithmic} regret,
$\mathrm{Reg}(T)=O\!\big(\sum_{k:\Delta_k>0}\tfrac{\log T}{\Delta_k}\big)$ with
$\Delta_k=\mu^\star-\mu_k$. \keyterm{Thompson sampling} reaches the same
logarithmic regret by a Bayesian route: keep a posterior over each arm's mean,
draw one sample from each, and pull the arm with the largest sample --- arms the
posterior is unsure about are sampled optimistically on their own. The left panel
of the figure shows both staying sub-linear while $\epsilon$-greedy's regret grows
as a straight line and eventually overtakes them.

% ======================================================================
\section{The policy gradient is the score function}
% ======================================================================

Now the harder problem: adjust a parameterised policy $\pi_\theta(a\mid s)$ to
maximise expected reward $J(\theta)=\mathbb{E}_{a\sim\pi_\theta}[R(a)]$. The reward
depends on $\theta$ only through the \emph{sampling} distribution, so we cannot
differentiate $R$ --- but we can differentiate the probability, with the
log-derivative trick of Chapter~2.

\begin{theorem}[Policy gradient / REINFORCE]
$\displaystyle \nabla_\theta J(\theta)
 = \mathbb{E}_{a\sim\pi_\theta}\!\big[R(a)\,\nabla_\theta\log\pi_\theta(a)\big].$
\end{theorem}
\begin{proof}
$\nabla_\theta J = \nabla_\theta\sum_a \pi_\theta(a)R(a)
 = \sum_a R(a)\,\nabla_\theta\pi_\theta(a)$. Multiply and divide by
$\pi_\theta(a)$ and use $\nabla_\theta\pi_\theta=\pi_\theta\nabla_\theta\log\pi_\theta$
(the log-derivative identity):
\[
  \nabla_\theta J = \sum_a \pi_\theta(a)\,R(a)\,\nabla_\theta\log\pi_\theta(a)
   = \mathbb{E}_{a\sim\pi_\theta}\big[R(a)\,\nabla_\theta\log\pi_\theta(a)\big].
\]
\end{proof}

The quantity $\nabla_\theta\log\pi_\theta(a)$ is the \keyterm{score} --- the exact
object from Chapter~2's maximum likelihood, now of a policy. The expectation is
estimated by sampling: take an action, observe its reward, and nudge $\theta$ in
the direction $R(a)\nabla\log\pi_\theta(a)$. This is \keyterm{REINFORCE}, and it
turns ``optimise through a sampler'' into an ordinary stochastic gradient ascent
(Chapter~3).

% ======================================================================
\section{Baselines are control variates}
% ======================================================================

REINFORCE is unbiased but noisy, because $R(a)$ can be large and its scale inflates
every gradient. The fix is a \keyterm{baseline} $b$: subtract it from the reward.

\begin{theorem}[A baseline does not bias the gradient]
For any constant $b$, $\;\mathbb{E}_{a\sim\pi_\theta}\!\big[(R(a)-b)\nabla\log\pi_\theta(a)\big]
= \nabla_\theta J(\theta)$.
\end{theorem}
\begin{proof}
It suffices that the subtracted term vanishes:
$\mathbb{E}[b\,\nabla\log\pi_\theta(a)] = b\sum_a \pi_\theta(a)\nabla\log\pi_\theta(a)
= b\sum_a \nabla\pi_\theta(a) = b\,\nabla\!\sum_a\pi_\theta(a) = b\,\nabla 1 = 0$,
using the mean-zero-score identity of Chapter~2.
\end{proof}

So the baseline is a \keyterm{control variate}: it changes the variance of the
estimator but not its mean, and the variance-minimising $b$ is a reward-magnitude
weighted average of the returns (close to, but not exactly, the mean reward). The
middle panel shows REINFORCE with a baseline learning faster and more smoothly
than without; the right panel plots the estimator's variance against $b$, a clean
parabola minimised near the mean reward. In full RL the natural baseline is the
\emph{value function} $V(s)$ (Chapter~6), and $R-V(s)$ is the \keyterm{advantage}
--- the actor--critic architecture is exactly this pairing. Finally, when the data
come from a different behaviour policy $\mu$, the same log-derivative machinery
corrects for it by \keyterm{importance sampling} (Chapter~4), reweighting each
sampled return by $\pi_\theta(a)/\mu(a)$ --- \keyterm{off-policy} policy
gradients.

\begin{summarybox}
In a $K$-armed bandit, regret is the price of not knowing the best arm; greedy
and $\epsilon$-greedy give linear regret, while UCB (an optimistic Hoeffding
confidence bound, Chapter~1) and Thompson sampling (a posterior draw) give
logarithmic regret. To optimise a sampled policy, the policy-gradient theorem
gives $\nabla_\theta J = \mathbb{E}[R(a)\nabla\log\pi_\theta(a)]$ via the
log-derivative trick, where $\nabla\log\pi_\theta$ is the Chapter~2 score;
REINFORCE is the stochastic-gradient estimate. A baseline $b$ leaves the gradient
unbiased (mean-zero score) but cuts its variance --- a control variate --- with
the value function as the natural choice ($R-V$ is the advantage, i.e.\
actor--critic), and importance sampling (Chapter~4) extends it off-policy.
\end{summarybox}

\end{document}
